\newcommand{\DOCMATIERE}{Chimie}\newcommand{\DOCNIVEAU}{Tle C D E}\newcommand{\DOCMODULE}{Module 2 · Acides et bases}\newcommand{\DOCLECON}{08}\newcommand{\DOCTITRE}{Acides faibles, bases faibles et couples acide/base}
% OnBuch+ — préambule « cours signés » ENRICHI (compilé avec tectonic / XeLaTeX)
% Système visuel « Pop » d'OnBuch+ : fond crème (jamais blanc), bordures encre
% épaisses, ombres « dures » décalées, titres Archivo Black en capitales.
% Version enrichie pour les cours de chimie : chimie (mhchem, chemfig),
% graphes (pgfplots), boîtes pédagogiques supplémentaires (définition,
% propriété, méthode, attention, le savais-tu, exercice résolu, exercice,
% TP, bilan), page de garde refaite et signature OnBuch+ ancrée en bas.
%
% Macros attendues dans main.tex AVANT \input{preamble} :
%   \DOCMATIERE  \DOCNIVEAU  \DOCMODULE  \DOCLECON (numéro)  \DOCTITRE
\documentclass[11pt]{article}

\usepackage{fontspec}
\usepackage[french]{babel}
\usepackage[a4paper,margin=2cm,top=3.4cm,bottom=2.8cm,headheight=1.4cm,footskip=1.3cm]{geometry}
\usepackage{amsmath,amssymb,amsfonts}
\usepackage[table]{xcolor} % [table] : fournit aussi \rowcolors (alternance de lignes)
\usepackage{pagecolor}
\usepackage{tikz}
\usetikzlibrary{arrows.meta,positioning,calc,shapes.geometric,shapes.misc,shapes.arrows,decorations.pathmorphing,decorations.markings,patterns,shadows,mindmap,trees,fit,backgrounds,matrix,intersections}
\usepackage{pgfplots}
\pgfplotsset{compat=1.18}
\usepackage[version=4]{mhchem}
\usepackage{chemfig}
\usepackage{enumitem}
\usepackage{fancyhdr}
% [most] charge toutes les bibliothèques tcolorbox (listings, minted, raster,
% documentation...) et gonfle inutilement la mémoire TeX sur un document long
% (~300 boîtes) : on ne charge que ce qui est réellement utilisé.
\usepackage[skins,breakable]{tcolorbox}
\usepackage{graphicx}
\usepackage{colortbl}
\usepackage{booktabs}
\usepackage{tabularx}
\usepackage{array}
\usepackage{multicol}
\usepackage{siunitx}
% parse-numbers=false : accepte aussi \SI{2,5 \times 10^{-3}}{...}
\sisetup{locale=FR,output-decimal-marker={,},group-minimum-digits=4,parse-numbers=false}
\DeclareSIUnit\ohm{\ensuremath{\symup{\Omega}}} % U+2126 absent de la police
\usepackage{qrcode}
% \degree : commande fréquemment utilisée par les modèles pour un angle ou une
% température en dehors de \SI{}{} (non fournie par les paquets chargés).
\providecommand{\degree}{^{\circ}}
\usepackage{lastpage}
\usepackage{hyperref}
\hypersetup{hidelinks}

% ============================================================
% Palette « Pop » OnBuch+ — mêmes hex que la classe P (plus_ui.dart)
% ============================================================
\definecolor{poporange}{HTML}{F59321}
\definecolor{popdark}{HTML}{E07A0C}
\definecolor{popcream}{HTML}{FFF6E8}
\definecolor{popink}{HTML}{1C1714}
\definecolor{poppurple}{HTML}{7A5AE0}
\definecolor{popgreen}{HTML}{2E9E6A}
\definecolor{poppink}{HTML}{E0457B}
\definecolor{popblue}{HTML}{2F7DE1}
\definecolor{popgold}{HTML}{C98A12}
\definecolor{popmuted}{HTML}{8A7F76}
\definecolor{popline}{HTML}{EADFD2}
\definecolor{popgray}{HTML}{8A7F76}
\colorlet{popgrey}{popgray}
% Alias tolérants (le modèle nomme parfois les couleurs différemment)
\colorlet{popwhite}{white}
\colorlet{popblack}{popink}
\colorlet{popred}{poppink}
\colorlet{popyellow}{popgold}
% Teintes claires pour fonds de tableaux / zones
\colorlet{popblueL}{popblue!10!white}
\colorlet{poporangeL}{poporange!14!white}
\colorlet{popgreenL}{popgreen!12!white}
\colorlet{poppurpleL}{poppurple!12!white}
\colorlet{poppinkL}{poppink!12!white}
\colorlet{popgoldL}{popgold!14!white}

\pagecolor{popcream}

% ============================================================
% Typographie
% ============================================================
\defaultfontfeatures{Ligatures=TeX}
\setmainfont{PlusJakartaSans-Medium.ttf}[Path=../../fonts/,BoldFont=PlusJakartaSans-Bold.ttf]
% Maths en sans-serif (Fira Math) avec chiffres et lettres droites de Plus
% Jakarta, pour que les formules se fondent dans le texte.
\usepackage[math-style=ISO,bold-style=ISO]{unicode-math}
\setmathfont{FiraMath-Regular.otf}[Path=../../fonts/]
\setmathfont{PlusJakartaSans-Medium.ttf}[Path=../../fonts/,range={up/num,up/{Latin,latin},\mathup}]
\usepackage{icomma} % virgule décimale sans espace en mode math
% Caractères absents de la police de texte (sinon carré vide) : rendus par
% leur équivalent mathématique, en texte comme en formule.
\usepackage{newunicodechar}
\newunicodechar{α}{\ensuremath{\alpha}}
\newunicodechar{Α}{\ensuremath{\alpha}}
\newunicodechar{β}{\ensuremath{\beta}}
\newunicodechar{δ}{\ensuremath{\delta}}
\newunicodechar{✓}{\popcheck{popgreen}}
\newfontfamily\popdisplay{ArchivoBlack-Regular.ttf}[Path=../../fonts/]
\newfontfamily\poplogo{SpaceGrotesk-Bold.ttf}[Path=../../fonts/]
\newfontfamily\popsemi{PlusJakartaSans-SemiBold.ttf}[Path=../../fonts/]
\newfontfamily\popxbold{PlusJakartaSans-ExtraBold.ttf}[Path=../../fonts/]
\newfontfamily\popxboldls{PlusJakartaSans-ExtraBold.ttf}[Path=../../fonts/,LetterSpace=3.0]

\setlist[enumerate]{leftmargin=1.7em,itemsep=5pt,topsep=4pt}
\setlist[itemize]{leftmargin=1.7em,itemsep=4pt,topsep=4pt,label={\color{poporange}\textbullet}}
\setlength{\parindent}{0pt}
\setlength{\parskip}{5pt}
\renewcommand{\arraystretch}{1.25}

% chemfig : liaisons un peu plus épaisses, encre Pop
\setchemfig{atom sep=2em,bond style={line width=0.9pt,popink},cram width=3pt}

% pgfplots : style Pop par défaut
\pgfplotsset{
  every axis/.append style={axis line style={popink,line width=1pt},tick style={popink},
    grid style={popline,line width=0.5pt},label style={font=\small},tick label style={font=\footnotesize}},
  popcurve/.style={poporange,line width=1.8pt,no markers,smooth},
}

% ============================================================
% Pill / squiggle
% ============================================================
\newcommand{\poppill}[3]{%
  \tikz[baseline=-0.55ex]\node[draw=popink,line width=1.2pt,rounded corners=2.6pt,%
    fill=#2,inner xsep=6.5pt,inner ysep=3pt]{%
    \popxboldls\fontsize{7}{7}\selectfont\color{#3}\MakeUppercase{#1}};%
}
\newcommand{\popsquiggle}[1]{%
  \tikz[baseline=-0.4ex]{
    \draw[#1,line width=2.6pt,line cap=round]
      plot [smooth] coordinates {(0,0.09) (0.34,0.01) (0.68,0.21) (1.02,0.01) (1.36,0.10)};
  }%
}
% Mot-clé surligné (marqueur orange sous le texte)
\newcommand{\cle}[1]{{\popxbold\color{popdark}#1}}

% ============================================================
% En-tête / pied
% ============================================================
\pagestyle{fancy}
\fancyhf{}
\renewcommand{\headrulewidth}{0pt}
\renewcommand{\footrulewidth}{0pt}
\lhead{\raisebox{-0.15ex}{\poplogo\fontsize{13}{13}\selectfont\color{popink}OnBuch\color{poporange}+}}
\rhead{\poppill{\DOCMATIERE\ \textbullet\ \DOCNIVEAU\ \textbullet\ Leçon \DOCLECON}{white}{popink}}
\lfoot{\popxbold\fontsize{7.6}{7.6}\selectfont\color{popmuted}\MakeUppercase{Cours signé OnBuch+}\ \textbullet\ {\popsemi\DOCTITRE}}
\rfoot{\tikz[baseline=-0.6ex]\node[rounded corners=8.5pt,fill=popink,minimum height=17pt,minimum width=17pt,inner xsep=5pt,inner sep=0]{\popxbold\fontsize{8}{8}\selectfont\color{popcream}\thepage\,/\,\pageref*{LastPage}};}

% ============================================================
% Titres
% ============================================================
\newcommand{\chaptitre}[2]{%
  \begin{tcolorbox}[enhanced,colback=poporange,colframe=popink,boxrule=2.6pt,arc=11pt,
      drop shadow=popink,left=15pt,right=15pt,top=12pt,bottom=11pt,before skip=3pt,after skip=18pt]
    \poppill{#2}{popink}{popcream}\par\vspace{9pt}
    {\raggedright\hyphenpenalty=10000\exhyphenpenalty=10000\popdisplay\fontsize{22}{24}\selectfont\color{popcream}\MakeUppercase{#1}\par}
  \end{tcolorbox}
}
% Section numérotée : pastille orange avec le numéro + titre Archivo
\newcounter{popsec}
\newcommand{\coursec}[1]{%
  \stepcounter{popsec}\setcounter{popsub}{0}%
  \par\vspace{16pt}\noindent
  \begin{minipage}{\linewidth}
  \tikz[baseline=-0.7ex]\node[draw=popink,line width=1.6pt,fill=poporange,rounded corners=4pt,minimum size=22pt,inner sep=0]{\popdisplay\fontsize{12}{12}\selectfont\color{popcream}\thepopsec};%
  \hspace{8pt}{\popdisplay\fontsize{15}{17}\selectfont\color{popink}\MakeUppercase{#1}}\par
  \vspace{4pt}\noindent{\color{poporange}\rule{36pt}{3.6pt}}
  \end{minipage}\par\nopagebreak\vspace{8pt}
}
\newcounter{popsub}
\newcommand{\courssub}[1]{%
  \stepcounter{popsub}%
  \par\vspace{9pt}\noindent{\popxbold\fontsize{11.5}{13}\selectfont\color{popink}{\color{poporange}\thepopsec.\thepopsub}\hspace{6pt}#1}\par\nopagebreak\vspace{4pt}
}

% ============================================================
% Boîtes pédagogiques — même recette : fond blanc, bordure épaisse colorée,
% ombre dure encre, badge plein. Titre optionnel [#1].
% ============================================================
\tcbset{popbase/.style={breakable,enhanced,colback=white,boxrule=2.3pt,arc=9pt,
  shadow={3pt}{-3pt}{0pt}{popink},left=14pt,right=14pt,top=11pt,bottom=11pt,before skip=11pt,after skip=15pt,
  fontupper=\popsemi\fontsize{10.6}{14.5}\selectfont\color{popink},
  fontlower=\popsemi\fontsize{10.6}{14.5}\selectfont\color{popink}}}
% Coche dessinée (le glyphe ✓ n'existe pas dans les polices du document)
\newcommand{\popcheck}[1]{\tikz[baseline=-0.5ex]\draw[#1,line width=1.6pt,line cap=round,line join=round](-0.13,0.0)--(-0.04,-0.09)--(0.13,0.1);}
\newcommand{\popboxhead}[3]{\poppill{#1}{#2}{white}\ifx\relax#3\relax\else\hspace{6pt}{\popxbold\fontsize{10.6}{12}\selectfont\color{popink}#3}\fi\par\vspace{7pt}}

\newtcolorbox{aretenir}[1][]{popbase,colframe=popgreen,before upper={\popboxhead{À retenir}{popgreen}{#1}}}
\newtcolorbox{exemplebox}[1][]{popbase,colframe=poporange,before upper={\popboxhead{Exemple}{poporange}{#1}}}
\newtcolorbox{definition}[1][]{popbase,colframe=popblue,before upper={\popboxhead{Définition}{popblue}{#1}}}
\newtcolorbox{propriete}[1][]{popbase,colframe=poppurple,before upper={\popboxhead{Propriété}{poppurple}{#1}}}
\newtcolorbox{methode}[1][]{popbase,colframe=popgold,before upper={\popboxhead{Méthode}{popgold}{#1}}}
\newtcolorbox{attention}[1][]{popbase,colframe=poppink,before upper={\popboxhead{Attention}{poppink}{#1}}}
\newtcolorbox{experience}[1][]{popbase,colframe=popblue,colback=popblueL,before upper={\popboxhead{Expérience}{popblue}{#1}}}
% « Le savais-tu ? » : carte violette pleine (contexte, culture, Cameroun)
\newtcolorbox{savaistu}[1][]{popbase,colback=poppurple,colframe=popink,
  fontupper=\popsemi\fontsize{10.4}{14.2}\selectfont\color{white},
  before upper={\poppill{Le savais-tu\,?}{white}{poppurple}\ifx\relax#1\relax\else\hspace{6pt}{\popxbold\color{white}#1}\fi\par\vspace{7pt}}}
% Exercice résolu : énoncé en haut, \tcblower puis la solution sur fond crème
\newcounter{popexo}\newcounter{popexr}
\newtcolorbox{exoresolu}[1][]{popbase,colframe=poporange,colbacklower=popcream,
  segmentation style={popink,line width=1.2pt,dashed},
  before upper={\stepcounter{popexr}\popboxhead{Exercice résolu \thepopexr}{poporange}{#1}},
  before lower={\poppill{Solution}{popink}{popcream}\par\vspace{6pt}}}
% Exercice (non résolu en place) — la correction est dans la partie Corrigés
\newtcolorbox{exercice}[1][]{popbase,colframe=popink,
  before upper={\stepcounter{popexo}\popboxhead{Exercice \thepopexo}{popink}{#1}}}
% Corrigé
\newtcolorbox{corrige}[1][]{popbase,colframe=popgreen,colback=popgreenL,
  before upper={\popboxhead{Corrigé}{popgreen}{#1}}}
% Objectifs / prérequis / bilan
\newtcolorbox{objectifs}{popbase,colframe=popink,colback=poporangeL,
  before upper={\popboxhead{Objectifs de la leçon}{popink}{}}}
\newtcolorbox{prerequis}{popbase,colframe=popink,colback=popblueL,
  before upper={\popboxhead{Prérequis}{popblue}{}}}
\newtcolorbox{bilan}{popbase,colframe=popink,colback=white,boxrule=2.8pt,
  before upper={\popboxhead{Fiche bilan}{poporange}{L'essentiel à connaître pour l'examen}}}
% Figure encadrée avec légende
\newtcolorbox{popfigure}[1][]{enhanced,breakable=false,colback=white,colframe=popline,boxrule=1.4pt,arc=8pt,
  left=10pt,right=10pt,top=10pt,bottom=8pt,before skip=10pt,after skip=12pt,halign=center,
  lower separated=false,halign lower=center,
  fontlower=\popsemi\fontsize{9}{11.5}\selectfont\color{popmuted},
  #1}
\newcommand{\legende}[1]{\par\vspace{6pt}{\popsemi\fontsize{9}{11.5}\selectfont\color{popmuted}#1}}

% ============================================================
% Signature OnBuch+
% ============================================================
\newcommand{\poplogoinline}[1]{{\poplogo\fontsize{#1}{#1}\selectfont\color{popink}OnBuch\color{poporange}+}}
\newcommand{\signatureonbuch}{%
  \begin{tcolorbox}[enhanced,colback=popink,colframe=popink,boxrule=2.2pt,arc=10pt,
      drop shadow=poporange,left=14pt,right=14pt,top=10pt,bottom=10pt,before skip=0pt,after skip=0pt]
    \tikz[baseline=-0.7ex]\node[circle,fill=popgreen,draw=popcream,line width=1.3pt,minimum size=22pt,inner sep=0]{\popcheck{white}};%
    \hspace{9pt}\begin{minipage}[c]{0.62\linewidth}
      {\popxbold\fontsize{12}{13}\selectfont\color{popcream}Signé\ \textbullet\ {\poplogo OnBuch\color{poporange}+}}\par\vspace{2pt}
      {\raggedright\popsemi\fontsize{8}{10.5}\selectfont\color{popline}\MakeUppercase{Équipe pédagogique OnBuch+}\\\MakeUppercase{Conforme au programme officiel MINESEC}\par}
    \end{minipage}\hfill
    {\poplogo\fontsize{22}{22}\selectfont\color{popcream}OnBuch\color{poporange}+}
  \end{tcolorbox}%
}

% ============================================================
% Page de garde
% #1 titre · #2 matière · #3 niveau · #4 module · #5 n° leçon · #6 durée ·
% #7 description / objectifs (texte ou itemize)
% ============================================================
\newcommand{\pagegarde}[7]{%
  \thispagestyle{empty}
  \newgeometry{a4paper,margin=1.7cm,top=1.6cm,bottom=1.4cm}%
  \begin{tikzpicture}[remember picture,overlay]
    \fill[poporange!18] ([xshift=-3.2cm,yshift=-3.6cm]current page.north east) circle (4.6cm);
    \fill[poppurple!14] ([xshift=2.4cm,yshift=7.6cm]current page.south west) circle (3.8cm);
  \end{tikzpicture}%
  {\poplogo\fontsize{30}{30}\selectfont\color{popink}OnBuch\color{poporange}+}\hfill
  \raisebox{0.6ex}{\poppill{Cours signé \textbullet\ Premium}{popink}{popcream}}\par
  \vspace{1pt}\popsquiggle{poppurple}\par
  \vspace{8pt}
  \begin{tcolorbox}[enhanced,colback=poporange,colframe=popink,boxrule=2.8pt,arc=13pt,
      drop shadow=popink,left=18pt,right=18pt,top=12pt,bottom=13pt,after skip=10pt]
    \poppill{#2 \textbullet\ #3}{popink}{popcream}\hspace{5pt}\poppill{#4}{white}{popink}\par\vspace{8pt}
    {\popdisplay\fontsize{14}{14}\selectfont\color{popink}LEÇON #5}\par\vspace{4pt}
    {\raggedright\hyphenpenalty=10000\exhyphenpenalty=10000\popdisplay\fontsize{25}{27}\selectfont\color{popcream}\MakeUppercase{#1}\par}
    \vspace{8pt}
    {\popxbold\fontsize{9.5}{9.5}\selectfont\color{popink}\MakeUppercase{Durée conseillée : #6}}
  \end{tcolorbox}
  \begin{tcolorbox}[enhanced,colback=white,colframe=popink,boxrule=2.2pt,arc=10pt,
      drop shadow=popink,left=16pt,right=16pt,top=10pt,bottom=10pt,after skip=8pt]
    \poppill{Dans cette leçon}{poppurple}{white}\par\vspace{5pt}
    \popsemi\fontsize{9.3}{12.6}\selectfont\color{popink}#7
  \end{tcolorbox}
  \noindent\begin{minipage}[t]{0.487\linewidth}
    \begin{tcolorbox}[enhanced,colback=popgreen,colframe=popink,boxrule=2.2pt,arc=10pt,
        drop shadow=popink,left=11pt,right=11pt,top=10pt,bottom=10pt,height=3.1cm,valign=center]
      \tcbox[colback=white,colframe=popink,boxrule=1.3pt,arc=5pt,boxsep=0pt,nobeforeafter,
        left=3pt,right=3pt,top=3pt,bottom=3pt,valign=center]{\qrcode[height=1.9cm]{https://play.google.com/store/apps/details?id=com.luvvix.onbuch&hl=fr}}%
      \hspace{8pt}\begin{minipage}[c]{0.5\linewidth}
        {\popdisplay\fontsize{11}{12.5}\selectfont\color{white}\MakeUppercase{Télécharge OnBuch}}\par\vspace{4pt}
        {\raggedright\popsemi\fontsize{8.4}{11}\selectfont\color{white}Gratuit \textbullet\ Google Play\\Tuteur IA, annales, concours, résultats\par}
      \end{minipage}
    \end{tcolorbox}
  \end{minipage}\hfill
  \begin{minipage}[t]{0.487\linewidth}
    \begin{tcolorbox}[enhanced,colback=poppurple,colframe=popink,boxrule=2.2pt,arc=10pt,
        drop shadow=popink,left=11pt,right=11pt,top=10pt,bottom=10pt,height=3.1cm,valign=center]
      \tikz[baseline=-0.7ex]\node[circle,fill=white,draw=popink,line width=1.3pt,minimum size=1.6cm,inner sep=0]{\popdisplay\fontsize{24}{24}\selectfont\color{poppurple}+};%
      \hspace{8pt}\begin{minipage}[c]{0.6\linewidth}
        {\popdisplay\fontsize{11}{12.5}\selectfont\color{white}\MakeUppercase{Passe à OnBuch+}}\par\vspace{4pt}
        {\raggedright\popsemi\fontsize{8.4}{11}\selectfont\color{white}Tous les cours signés, Tuteur IA illimité, fiches PDF — onglet OnBuch+ de l'app\par}
      \end{minipage}
    \end{tcolorbox}
  \end{minipage}
  \vspace{8pt}
  \popcontenu
  \vfill
  \signatureonbuch
  \restoregeometry
  \newpage
}

% Rangée « Ce cours contient » (page de garde)
\newcommand{\poptile}[3]{% #1 couleur #2 gros texte #3 légende
  \begin{tcolorbox}[enhanced,colback=#1,colframe=popink,boxrule=2pt,arc=9pt,shadow={2.5pt}{-2.5pt}{0pt}{popink},
      left=6pt,right=6pt,top=9pt,bottom=9pt,halign=center,valign=center,height=1.75cm,width=\linewidth,nobeforeafter]
    {\popdisplay\fontsize{12.5}{13}\selectfont\color{white}\MakeUppercase{#2}}\par\vspace{3pt}
    {\centering\popsemi\fontsize{7.8}{9.5}\selectfont\color{white}#3\par}
  \end{tcolorbox}}
\newcommand{\popcontenu}{%
  \par\noindent{\popxboldls\fontsize{7.5}{7.5}\selectfont\color{popmuted}CE COURS SIGNÉ CONTIENT}\par\vspace{7pt}
  \noindent\begin{minipage}[t]{0.235\linewidth}\poptile{popblue}{Cours}{complet, illustré, pas à pas}\end{minipage}\hfill
  \begin{minipage}[t]{0.235\linewidth}\poptile{poporange}{Méthodes}{et exercices résolus}\end{minipage}\hfill
  \begin{minipage}[t]{0.235\linewidth}\poptile{poppink}{Exercices}{type examen + corrigés}\end{minipage}\hfill
  \begin{minipage}[t]{0.235\linewidth}\poptile{popgreen}{Fiche bilan}{l'essentiel à retenir}\end{minipage}\par}

% Sommaire
\newtcolorbox{sommaire}{enhanced,colback=white,colframe=popink,boxrule=2.2pt,arc=10pt,
  drop shadow=popink,left=18pt,right=18pt,top=13pt,bottom=13pt,before skip=6pt,after skip=20pt,
  before upper={\poppill{Sommaire}{poppurple}{white}\par\vspace{9pt}\popsemi\fontsize{10.6}{14.5}\selectfont\color{popink}}}

% Signature de fin de cours — poussée tout en bas de la dernière page
\newcommand{\findecours}{%
  \par\vfill\nopagebreak
  \begin{center}\popsquiggle{poporange}\end{center}\vspace{4pt}
  \signatureonbuch
}
% Glyphes absents de Fira Math : redéfinis à partir de glyphes disponibles
\AtBeginDocument{%
  \def\setminus{\mathbin{\backslash}}\let\smallsetminus\setminus
  \def\vdots{\mathord{\vbox{\baselineskip3.2pt\lineskiplimit0pt\kern4pt\hbox{.}\hbox{.}\hbox{.}}}}%
  \def\ddots{\mathinner{\mkern1mu\raise6.5pt\hbox{.}\mkern2mu\raise3.6pt\hbox{.}\mkern2mu\raise0.7pt\hbox{.}\mkern1mu}}%
  \def\triangle{\mathord{\scalebox{0.9}{$\symup{\Delta}$}}}%
  \def\nabla{\mathord{\raisebox{\depth}{\scalebox{1}[-1]{$\symup{\Delta}$}}}}%
  \def\checkmark{\popcheck{popdark}}%
  \def\star{\mathbin{\ast}}\def\bigstar{\mathord{\ast}}%
  \def\diamond{\mathbin{\scalebox{0.6}{$\Diamond$}}}%
  \def\bigcup{\mathop{\mathchoice{\raisebox{-0.3em}{\scalebox{1.7}{$\cup$}}}{\scalebox{1.25}{$\cup$}}{\cup}{\cup}}\displaylimits}%
  \def\bigcap{\mathop{\mathchoice{\raisebox{-0.3em}{\scalebox{1.7}{$\cap$}}}{\scalebox{1.25}{$\cap$}}{\cap}{\cap}}\displaylimits}%
  \def\lesseqqgtr{\mathrel{\lessgtr}}\def\bigoplus{\mathop{\oplus}\displaylimits}%
}
\providecommand{\ssi}{si et seulement si} % abréviation fréquente

\begin{document}
\begin{tcolorbox}[enhanced,colback=popink,colframe=popink,arc=10pt,boxrule=0pt,left=12pt,right=12pt,top=9pt,bottom=9pt]
{\color{white}\popxbold\small FICHE DE TD 3/3 \textbullet\ DIFFICILE \hfill Chimie \textbullet\ Tle C D E\par}
\vspace{4pt}{\color{white}\popxbold\Large Acides faibles, bases faibles et couples acide/base\par}
\vspace{3pt}{\color{white}\small Module 2 · Acides et bases\par}
\end{tcolorbox}
\vspace{4pt}
\noindent \textbf{Consigne :} Traiter les exercices avec rigueur, justifier chaque étape par les relations du cours et discuter les résultats obtenus.

\begin{exercice}[Titrage d'un acide faible par une base forte : étude complète]
On réalise le titrage au dosage de \SI{20,0}{\milli\liter} d'une solution aqueuse d'acide éthanoïque \ce{CH3COOH} (noté \ce{AH}) de concentration molaire inconnue $C_{\ce{AH}}$, par une solution aqueuse de soude \ce{NaOH} de concentration molaire $C_{\ce{B}} = \SI{0,100}{\mole\per\liter}$. On note \ce{A-} l'ion éthanoate. Le suivi pH-métrique donne les points caractéristiques suivants : 
\begin{itemize}
    \item Point initial : $\mathrm{pH}_0 = 2,88$.
    \item Point d'équivalence : $V_{\mathrm{eq}} = \SI{22,0}{\milli\liter}$, $\mathrm{pH}_{\mathrm{eq}} = 8,72$.
    \item Point demi-équivalence : $V_{\mathrm{eq}}/2 = \SI{11,0}{\milli\liter}$, $\mathrm{pH}_{1/2} = 4,76$.
\end{itemize}
$K_{\mathrm{e}} = \num{1,0e-14}$ à la température de l'expérience.

\begin{enumerate}
    \item \textbf{État initial.} À partir de $\mathrm{pH}_0$, calculer la concentration molaire $C_{\ce{AH}}$ de la solution titrée. En déduire le coefficient d'ionisation $\alpha_0$ et la constante d'acidité $K_{\mathrm{a}}$ du couple \ce{AH/A-}. Commenter la valeur de $\alpha_0$ au regard de la nature de l'acide.
    \item \textbf{Équivalence.} Écrire l'équation de la réaction de titrage et justifier que le pH à l'équivalence est basique. Calculer la concentration de l'ion \ce{A-} à l'équivalence. En déduire la constante de basicité $K_{\mathrm{b}}$ du couple \ce{A-/AH} et vérifier la relation $K_{\mathrm{a}} \times K_{\mathrm{b}} = K_{\mathrm{e}}$.
    \item \textbf{Domaine tampon.} Montrer que la relation $\mathrm{pH} = \mathrm{p}K_{\mathrm{a}} + \log\left(\frac{[\ce{A-}]}{[\ce{AH}]}\right)$ est vérifiée au point demi-équivalence. Déterminer l'intervalle de pH pour lequel le mélange \ce{AH/A-} constitue une solution tampon efficace. Justifier.
    \item \textbf{Choix de l'indicateur.} Parmi les indicateurs suivants, lequel choisir pour détecter l'équivalence visuellement ? Justifier par le domaine de virage et le saut de pH à l'équivalence.
    \begin{center}
    \begin{tabular}{lcc}
    \hline
    Indicateur & Domaine de virage (pH) & Couleur (acide $\to$ base) \\
    \hline
    Méthylorange & 3,1 -- 4,4 & Rouge $\to$ Jaune \\
    Bleu de bromothymol & 6,0 -- 7,6 & Jaune $\to$ Bleu \\
    Phénolphtaléine & 8,2 -- 10,0 & Incolore $\to$ Rose \\
    \hline
    \end{tabular}
    \end{center}
\end{enumerate}
\end{exercice}

\begin{exercice}[Hydrolyse d'un sel d'acide faible et de base forte : loi d'Ostwald]
On introduit $n_0 = \SI{5,0e-3}{\mole}$ d'éthanoate de sodium \ce{CH3COONa} (sel totalement dissocié) dans un ballon jaugé de \SI{250,0}{\milli\liter} que l'on complète à l'eau distillée. On note $C_0$ la concentration molaire formelle en ions \ce{A-}. Le pH de la solution finale est mesuré : $\mathrm{pH} = 8,95$.

\begin{enumerate}
    \item \textbf{Équilibre d'hydrolyse.} Écrire l'équation de la réaction d'hydrolyse de l'ion \ce{A-} (réaction avec l'eau). Préciser le couple acide/base mis en jeu. Exprimer l'avancement $x$ de cette réaction à l'équilibre en fonction de $[\ce{OH-}]_{\mathrm{eq}}$.
    \item \textbf{Constante de basicité.} Calculer $K_{\mathrm{b}}$ du couple \ce{A-/AH} à partir du pH mesuré. En déduire $K_{\mathrm{a}}$ du couple \ce{AH/A-} et $\mathrm{p}K_{\mathrm{a}}$. Comparer cette valeur à celle trouvée à l'exercice précédent (4,76) et discuter l'éventuel écart.
    \item \textbf{Coefficient d'hydrolyse et dilution.} Définir le coefficient d'hydrolyse $h$ (ou taux d'avancement final $\tau$). Calculer sa valeur ici. 
    On dilue ensuite la solution mère par 10 (prélèvement de \SI{25,0}{\milli\liter} complété à \SI{250,0}{\milli\liter}). Sans refaire tous les calculs, prévoir qualitativement l'évolution de $h$ et du pH. Justifier quantitativement à l'aide de la loi d'Ostwald (relation entre $h$, $K_{\mathrm{b}}$ et $C_0$) si l'approximation $h \ll 1$ reste valide.
    \item \textbf{Force relative.} Classer les couples \ce{H3O+/H2O}, \ce{AH/A-}, \ce{H2O/OH-}, \ce{A-/AH} par ordre de force acide croissante. En déduire l'ordre de force basique croissante des bases correspondantes. Expliquer pourquoi l'ion \ce{A-} est une base plus forte que l'eau.
\end{enumerate}
\end{exercice}

\begin{exercice}[Réactions entre couples acide/base : prévision et constante de réaction]
On considère les couples acide/base suivants avec leurs $\mathrm{p}K_{\mathrm{a}}$ (eau, \SI{25}{\celsius}) :
\begin{center}
\begin{tabular}{lcc}
\hline
Couple & $\mathrm{p}K_{\mathrm{a}}$ & Nature \\
\hline
\ce{NH4+/NH3} & 9,2 & Acide ammonium / Base ammoniac \\
\ce{CH3COOH/CH3COO-} & 4,76 & Acide éthanoïque / Base éthanoate \\
\ce{HCN/CN-} & 9,2 & Acide cyanhydrique / Base cyanure \\
\ce{HCO3-/CO3^{2-}} & 10,3 & Acide hydrogénocarbonate / Base carbonate \\
\hline
\end{tabular}
\end{center}

\begin{enumerate}
    \item \textbf{Classification.} Classer les acides \ce{NH4+}, \ce{CH3COOH}, \ce{HCN}, \ce{HCO3-} par ordre de force acide croissante. Classer les bases \ce{NH3}, \ce{CH3COO-}, \ce{CN-}, \ce{CO3^{2-}} par ordre de force basique croissante. Justifier à chaque fois par la valeur de $\mathrm{p}K_{\mathrm{a}}$.
    \item \textbf{Prévision de réaction.} On mélange des quantités de matière égales d'ions \ce{CH3COO-} et d'ions \ce{NH4+}. 
    \begin{enumerate}
        \item Écrire l'équation de la réaction éventuelle entre ces deux espèces. Identifier l'acide et la base réactifs.
        \item Calculer la constante de réaction $K$ associée. Conclure sur le sens d'évolution spontané de la réaction (totale, partielle, inexistante).
        \item Calculer le taux d'avancement final $\tau$ si les concentrations initiales sont $[\ce{CH3COO-}]_0 = [\ce{NH4+}]_0 = \SI{0,10}{\mole\per\liter}$.
    \end{enumerate}
    \item \textbf{Ampholytes et domaines de prédominance.} L'ion \ce{HCO3-} est un ampholyte. 
    \begin{enumerate}
        \item Écrire les deux équilibres acide/base dont il est l'acteur principal. Préciser les couples correspondants.
        \item Déterminer le domaine de pH pour lequel \ce{HCO3-} est l'espèce prédominante par rapport à \ce{H2CO3} et par rapport à \ce{CO3^{2-}}. Représenter schématiquement (diagramme de prédominance) les domaines de prédominance des trois espèces \ce{H2CO3}, \ce{HCO3-}, \ce{CO3^{2-}} en fonction du pH.
        \item On prépare une solution de \ce{NaHCO3} à \SI{0,01}{\mole\per\liter}. Estimer le pH de cette solution (formule approchée pour les ampholytes). Justifier la validité de l'approximation.
    \end{enumerate}
    \item \textbf{Synthèse critique.} Un élève affirme : « Puisque $\mathrm{p}K_{\mathrm{a}}(\ce{NH4+/NH3}) = \mathrm{p}K_{\mathrm{a}}(\ce{HCN/CN-}) = 9,2$, les couples \ce{NH4+/NH3} et \ce{HCN/CN-} ont exactement la même force acide, donc mélanger \ce{NH3} et \ce{CN-} ne provoquera aucune réaction. » Critiquer cette affirmation en utilisant la relation $K = 10^{\mathrm{p}K_{\mathrm{a2}} - \mathrm{p}K_{\mathrm{a1}}}$.
\end{enumerate}
\end{exercice}

\begin{exercice}[Détermination expérimentale de $\mathrm{p}K_{\mathrm{a}}$ par mesure de conductivité]
On souhaite déterminer la constante d'acidité $K_{\mathrm{a}}$ de l'acide éthanoïque par conductimétrie. On mesure la conductivité $\sigma$ de solutions d'acide éthanoïque de concentrations formelles $C$ variées. On note $\Lambda$ la conductivité molaire de la solution et $\Lambda^{\circ}$ la conductivité molaire à dilution infinie (valeur tabulée : $\Lambda^{\circ}_{\ce{CH3COOH}} = \SI{390,7}{\s\centi\meter\squared\per\mole}$ à \SI{25}{\celsius}).

\begin{enumerate}
    \item \textbf{Relation conductivité - avancement.} Rappeler la relation entre le coefficient d'ionisation $\alpha$ (ou taux d'avancement $\tau$), la conductivité molaire $\Lambda$ et $\Lambda^{\circ}$. Justifier cette relation.
    \item \textbf{Expression de $K_{\mathrm{a}}$ en fonction de $\Lambda$.} En supposant que l'approximation $1-\alpha \approx 1$ n'est \textbf{pas} faite, établir l'expression littérale de $K_{\mathrm{a}}$ en fonction de $C$, $\Lambda$ et $\Lambda^{\circ}$.
    \item \textbf{Méthode graphique (Ostwald).} Montrer qu'on peut écrire : $\frac{C \Lambda}{\Lambda^{\circ}(\Lambda^{\circ} - \Lambda)} = \frac{1}{K_{\mathrm{a}}} \cdot \frac{\Lambda}{\Lambda^{\circ}} + \frac{1}{K_{\mathrm{a}}}$. Préciser l'ordonnée à l'origine et la pente du droit obtenu en représentant $y = \frac{C \Lambda}{\Lambda^{\circ}(\Lambda^{\circ} - \Lambda)}$ en fonction de $x = \frac{\Lambda}{\Lambda^{\circ}}$. Comment déterminer $K_{\mathrm{a}}$ graphiquement ?
    \item \textbf{Analyse critique.} Pourquoi la méthode conductimétrique est-elle particulièrement adaptée aux acides faibles ? Quelle serait la difficulté majeure si l'on tentait d'appliquer cette même méthode à un acide fort comme \ce{HCl} ? Discuter de l'influence de la température sur les mesures.
\end{enumerate}
\end{exercice}
\end{document}
